0\chapter{T\'itulo do primeiro cap\'itulo} \label{cap:cap1}

Escrever bem é uma arte que exige muita técnica e dedicação. Há vários bons livros sobre como escrever uma boa dissertação ou tese. Para a escrita de textos em Ciência da Computação, Writing for Computer Science \cite{zobel2014} é uma leitura obrigatória. O livro Metodologia de Pesquisa para Ciência da Computação \cite{wazlawick2009} também merece uma boa lida. 

\section{Introdu\c{c}\~ao}\label{sec:intro}

Apresentar uma visão geral do assunto que será abordado no trabalho, procurando fazer com que o leitor adquira uma compreensão inicial do que será tratado e fornecendo informações que o levem a perceber a sua importância.  

Exemplo de Figura: Ver Figura de Exemplo~\ref{fig:exefig}.

\begin{figure}[!ht]
\centering
\includegraphics[width=0.3\linewidth]{figuras/exefig.eps}
\caption{Exemplo de figura}
\label{fig:exefig}
\end{figure}


\section{Objetivos}\label{sec:intro}

O objetivo da pesquisa deve ser diretamente verificável ao final do trabalho. Um bom objetivo de pesquisa possivelmente irá demonstrar que alguma hipótese sendo testada é ou não verdadeira.
Portanto, o objetivo geral e os objetivos específicos do trabalho devem ser expressos na forma de uma condição não trivial cujo sucesso possa vir a ser verificado ao final do trabalho. Um objetivo bem expresso em geral terá verbos como “demonstrar”, “provar”, “melhorar” (de acordo com alguma métrica definida) etc.
Deve-se tomar cuidado com certos verbos que determinam objetivos cuja verificação é trivial e, portanto, inadequada. Entre eles pode-se citar “propor”, “estudar”, “apresentar” etc. Se o objetivo do trabalho é propor algo, basta que a coisa seja proposta para que o objetivo seja atingido e, portanto, essa forma é trivial e inadequada, pois a definição do objetivo não menciona a qualidade daquilo que será proposto.
Se o objetivo do trabalho é estudar algo, então ele terá sido alcançado se aquilo foi estudado, não importando se alguma nova informação foi aprendida ou não, sendo, portanto, inadequado como objetivo de pesquisa. Estudar, normalmente, é o objetivo do aluno e não do trabalho.
Se o objetivo do trabalho consiste em apresentar algo, novamente ele é trivial e inadequado. Uma simples apresentação não produz necessariamente conhecimento novo. Por exemplo, “o objetivo deste trabalho é apresentar os
operadores da lógica booleana”; tal objetivo pode ser alcançado com um pequeno texto explicando os operadores conhecidos, mas, como não traz informação nova, não é um objetivo de pesquisa.
A proposta, o estudo e a apresentação podem ser justificáveis como objetivo de pesquisa desde que o objeto da proposta, estudo ou da apresentação seja algo original.
\cite{wazlawick2009}.
Insira os objetivos específicos no formato de tópicos.

Exemplo de refer\^encias~\cite{tese2011, confinter2011, confnac2011, rev2011, site2011}.

Este \'e o primeiro cap\'itulo. 

Exemplo de Tabela: ver Tabela~\ref{tab:exetab}.

\begin{table}[!ht]
\begin{center}
\caption{Distribuição IMC em Adultos}
\label{tab:exetab}
\begin{tabular}{|c |c |}
\hline
\textbf{\textbf{Classificação}} & \textbf{IMC}\\
\hline\hline
Baixo Peso & < 18,5 \\
Peso Adequado & > 18,5 e < 25 \\
Sobrepeso & > 25 e < 30 \\
Obesidade & > 30 \\
\hline
\end{tabular}
\end{center}
\end{table}

Este \'e o primeiro cap\'itulo. 
 


\section{Desenvolvimento}\label{sec:desenv}

Conforme explicado na Se\c{c}\~ao~\ref{sec:intro}, nesta segunda se\c{c}\~ao ser\'a descrito o desenvolvimento.
